\documentclass[../../main.tex]{subfiles}
\graphicspath{{\subfix{../../image/}}}

\begin{document}

\section{注释与参考文献}\label{section:ch09-8-参考文献}

\paragraph{\S9.1。}见 J.-L. Lions [L2] 与 Zeidler [Zd, 第 2 卷]。小册子 [E1] 是关于处理非线性 PDE 弱收敛问题方法的综述。

\paragraph{\S9.2。}Zeidler [Zd, 第 1 卷] 有更多关于不动点技巧的内容。另见 Gilbarg--Trudinger [G-T, 第 11 章]。S. Armstrong 协助了 Schaefer 定理。

\paragraph{\S9.3。}见 Smoller [S, 第 10 章]。

\paragraph{\S9.4。}\S9.4.1 基于 Payne [Pa, \S9]。\S9.4.1\cref{theorem:ch09-4-2-小初值爆破} 归功于 H. Fujita（J.\ Fac.\ Sci.\ U.\ Tokyo Sect.\ I 13 (1966), 109--124）。

\paragraph{\S9.5。}\S9.5.2 归功于 Gidas--Ni--Nirenberg（Comm.\ Math.\ Phys.\ 68 (1979), 209--243）。

\paragraph{\S9.6。}这部分材料来自 Brezis [Br2]，该书包含更多关于 Hilbert 空间上非线性半群的内容。Banach 空间上的非线性半群理论见 Barbu [Ba] 或 Zeidler [Zd, 第 2 卷]。

\end{document}
